\section{Spiritual Pilgrims: Carl Jung and Teresa of Avila}

Father John Welch is a member of the Carmelite Order and a Jungian who has lectured and written extensively on the meaning and practice of \textbf{St. Teresa d'Avila}'s \emph{Interior Castle} in tandem with Jungian psychology, especially where it concerns the process of Individuation. Father Welch's lectures and writings serve equally as a primer to both subjects, or fresh, insightful reviews, for those already familiar with the territory. In various contemporary ``traditional'' writings and circles, both ``mysticism'' and Jungianism have often gotten short-shrift, or been subject to certain austere, and or carte blanche criticisms---thus the work of Father Welch serves to help counter balance these equations.

In \emph{Meditations on the Tarot}, Tomberg raises a few points worth considering in the above context:

\begin{enumerate}


\item The objective of this work is a ``fusion'' of types of knowledge.~ Hermetism seeks a resolution of the ``oppositions'' between spirituality and intellectualism, toward an agreeable ``parallelism'', and lastly a ``fusion''.~ This fusion Tomberg equates with the ``Philosopher's Stone'', and ``Operation of the Sun and the Moon'' (Meditation on The Fool).

\item That the work of Carl Jung constitutes a branch of such ``fusion'' in the domain of psychological and behavior sciences, Tomberg adds

\begin{quotex}
Neither \textbf{Vladimir Solovyov}, nor \textbf{Nicholas Berdyaev} nor \textbf{Pierre Teilhard de Chardin}, nor \textbf{Carl Gustav Jung}, for example were declared Hermeticists, but how much they have contributed to the progress of the work in question! Christian existentialism (Berdyaev), Christian Gnosis (Solovyov), Christian evolutionism (Chardin), and depth psychology of revelation (Jung) are, in fact, as many inestimable contributions to the cause of the union of spirituality and intellectuality 
\flright{\textit{Meditation on The Fool}}
\end{quotex}


In the specifically religious domain, Tomberg finds this type of fusion present where prayer and meditation coexist as method, or exercise. He offers as examples St. Teresa and St. John of the Cross in Carmelitism, St. Boneventura and the Franciscans, and St. Ignatius Loyola and the Jesuit Order (Meditation on The Fool).

\item That the beginning task of the spiritual aspirant, the ``initiate'' need attend to, by whatever other name one wants to describe it, what Jung dubbed ``individuation''.~ Tomberg explains that this individuation is the ``harmonization'' of the consciousness and unconsciousness ---this has also been referred to in the Hermetic literature as the ``chemical wedding'', preceding the ``mystical wedding''---in St. Teresa, as we will see, the individuation stage occurs throughout several ``rooms'', and culminates in ``betrothal''; whereas the seventh and final ``room'' corresponds to the ``mystical marriage''.~ In Tomberg, St. Teresa and Jung, the common result of this phase of development is the birth of a new ``Self'' which becomes the center of the being.
\end{enumerate}


\paragraph{Architecture of The Interior Castle}


St. Teresa depicts seven ``rooms'' in her castle, through which the athlete experiences differing intensities and effects of prayer, as well as numerous spiritual and external life challenges. Those familiar with Sufism might detect similarities between the castle chambers, with the Sufistic ``stations'' and their associated ``states''. In any event, the challenges of St. Teresa's path toward Union also evidence again, that unlike the shallow, ``fluffy bunny'', ``let's all think positive and attract good vibes from the universe'', pop-spirituality, here in the presence of an experienced heavy weight champion, we are taught that traveling the path may often exasperate internal and external difficulties, long before remedying such. In certain rooms, Teresa makes no bones about pointing out that one might feel alienated from both God and man, isolated, depressed, and longing for death. This too is the message of Jung concerning individuation, which he suggests is one of the most painful processes a human might undergo---and which is also why he avers, many adults quite simply abandon the process early on in their adult lives, never ``growing'' much beyond who they were in their twenties/early thirties. Keeping this last point in mind, when out in public next time, try unobtrusively observing and listening to strangers noting their age groups---this makes for an interesting thought experiment in itself.

Now despite the use of ``seven'' concrete rooms, one ought regard these as a matter of literary convenience, as the actual ``number'' of rooms is indefinite. Perhaps while considering the Gospel saying that ``In my Father's house there are many mansions'', St. Teresa noted that the castle is arranged in concentric rings, with each main room, in fact enveloping numerous sub-rooms, reminiscent of a palm frond as she described them---or maybe in these days we could call to mind imagery of fractals or Manderblots:

\begin{enumerate}


\item Room one is straight forward enough. The individual meets the most basic religious requirements---they ``pray'', but their prayer is on the cold side. In terms of lifestyle, they are preoccupied with most everything external as their priorities---in fact, although ``religious'', they have little awareness of, or interest in anything ``internal''. This room is crawling with ``reptiles'' and vermin---a menagerie of Jungian ``shadow'' symbols.

\item The second room is reached when the person begins to have an inkling that their life has a purpose, call, or ``vocation'', and proper discernment means developing a more intimate exchange with God. In prayer life, the person is becoming warmer, and they are exploring a bit more of scripture, liturgy, and charitable works.

\item This third room is particularly interesting, because it characterizes much of the ``exoterism'' of our times. But, the good news first. In the third room, the individual is living a ``well ordered life''. They are ``mature'' in their familial, social, professional, and religious lives---they basically have mastered the ``letter of the law'', and conduct themselves accordingly---``good Christians'', and what would be considered secularly as ``upstanding citizens''. Teresa explains though, that remaining too long in this station of sincere, but exclusively outward prayer results in spiritual stasis, and dryness. In a Jungian context, this station might roughly run from a person's mid-twenties to mid-thrities; while necessary, useful, and inherently positive, we might observe that when becoming irrevocably habitual, not only seizes up the growth process, but generates what might be seen as a ``bourgeois'' moralizing, and at the worse, a burgeoning fundamentalism closing itself off from the ``spirit of the law''.

\item Our fourth room is of great importance, as it establishes a division of the rooms into two groups, two broad hemispheres, representing the phases of individuation, with this fourth as ``transition''. From a Jungian perspective, the individuation process can be divided into two distinct phases:

\begin{enumerate}
  
\item The first half of life (approaching ages 35-40), marked by personality growth, and adaptation to the external environment, ``the world''.
  
\item Second ``half'' of life (beyond ages 35), and relating to the interior world.~ It should be noted that these ``ages'' are relatively relative, as Jung observed that some teen-aged patients had already embarked along the second half; while many adults well into advanced age had never even begun such as stage.
\end{enumerate}
\end{enumerate}


Ideally, the second phase produces an ``equilibrium'' between consciousness, and the unconscious poles within the psyche; the arrival at individuation, is the arrival at the ``Self''. Not for nothing then, did Teresa explain that while self-knowledge is unnecessary to obtain salvation, pursuit of increasingly more profound self-knowledge develops and promotes union with God. Thus, we are reminded of something Guenon and others have mentioned: exotericism, and the fulfillment of the rites of a religion, are available to all (the masses) to obtain salvation; while, others upon accomplishing this task, might seek ``more'', an esoterism, attempting to move the Self into the occupation of the higher states.

Father Welch has then approached the rooms of the Interior Castle as spiritual pilgrimage, as well as equally corresponding to the individuation phases, where rooms 1-3 correspond to Jungian phase I, rooms 5-7 Jungian phase II, and room 4, as transition pivot. Both Teresa and Jung agree that room 4, or the transition between phases I. and II. is a particularly difficult and trying station. In Hermetic language, this room/station reminds us of the ``mastery of the four elements'', or the ``trials of the Sphinx''. In this room, opposite pairs from within the four ``functional types'' (thinking, feeling, intuiting, sensing) emerge to hopefully enter equilibrium. A recently published book attributed to \textbf{Franz Bardon}, entitled \emph{The Universal Master Key} examines these functions or properties of the elements in the soul, from a Hermetic angle, with advice on developing them in harmony, and worth exploring in room four.

As pertains to prayer, Teresa's fourth room inaugurates the ``fusion'' or ``harmonization'' of outward prayer, and inner contemplation; it recalls what Tomberg mentioned above concerning intellectuality and spirituality.

5, 6, 7. Likely as one would anticipate then, the fifth, sixth, and seventh rooms, are ``hot'', individuation has reached an apex, and the spiritual life is characterized increasingly by states of Divine Union---the influence of the merciful descent of Grace tends to take greater precedence though than individual actions or ``exercises'' (though these are not dispensed with).

In his meditation on ``The Pope'', Tomberg discusses the meaning of ``benediction'', what is also called ``barakah'' in other traditions, and assigned the nomenclature of ``spiritual respiration'' in Tomberg. It involves another type of ``harmonizing''; Tomberg's imagery uses the left and right columns of the Kabbalistic Tree, the Boaz and Jachin pillars---more to the point, the upward ascent of prayer, followed by descent of benediction---never forgetting the critical matter, that this descent is a gift of mercy---not a product of a mechanical and ``scientific'' operation of ``arbitrary magic''. Two scriptural verses and a liturgical example help illustrate these principles:

\begin{quotex}
Many are called, but few are chosen.
\flright{Matthew 22:14}

Let my prayer ascend to you like incense, and the lifting up of my hands as an evening sacrifice.
\flright{Psalm 141:2}
\end{quotex}

And, the Epiclesis of the Mass/Divine Liturgy.

St. Teresa's language regarding the sixth room involves betrothal, and the seventh the ``Spiritual Marriage'', the reaching of the central room of the castle. For Jung, the center of the ``mandala'' has been reached; the whole Self attained. ``Horizontally'' the athlete, working in accord with the external law weds and makes whole the psyche mostly through their own efforts, energies, and disciplines, and a new ``center'' takes over. ``Vertically'', and more so through the grace of benediction, the harmonized soul might ascend to enter the bridal chamber of Spirit.

\begin{quotex}
And I, if I am lifted up from the earth, will draw all men unto me
\flright{John 12:23}
\end{quotex}

Notes:

\textbf{Jung, Carl}, \textit{Memories, Dreams, and Reflections, Undiscovered Self}.

\textbf{Teresa of Avila}, \textit{The Interior Castle}.\footnote{Online pdf at \url{http://www.saintsbooks.net}}

\textbf{Tomberg, Valentin}, \textit{Meditations on the Tarot}.


\textbf{Welch, John}, \textit{Spiritual Pilgrims: Carl Jung and Teresa of Avila}. Video lecture \textit{The Interior Castle of Teresa Avila: A Map of our Spiritual Journey} in Youtube.


\flrightit{Posted on 2014-04-03 by Frater M}
